\chapter{Introduction}
\label{ch:introduction}


This chapter states the problem, defines the term the report adjudicates, sets out the aim
and objectives, lists the contributions, explains the change of direction since the interim
report, and states the ethical boundary the work held.

\section{Motivation}
\label{sec:motivation}

A prediction market pays people for being right about the future. To do that it must do two
things that look administrative and are not. It must take a participant's belief and turn it
into a position, and it must later decide what actually happened. Both steps require the
belief to be exposed: once to whoever prices the position, and once to whoever judges the
outcome.

In a market where the participants and the resolver are human, those exposures are familiar
and their costs are understood. A resting order on a public book tells other traders
something; a resolution committee interprets a question and can be argued with. A new
generation of prediction markets changes both ends at once. Forecasting agents submit
positions programmatically, market makers are automated and always watching, and the party
deciding what happened is increasingly a language model executing a procedure that somebody
wrote in advance. Two protocols now running make opposite choices at both points. Delphi,
built on Gensyn's infrastructure, matches against a pari-mutuel pool and settles by AI
judgement. Sapience matches by broadcasting a request for quotes and resolves through a
bonded optimistic oracle.

This project asks a question that follows directly: when a market's traded object is a
belief and its resolver is a model, where does strategic informational advantage sit, what
is it worth, and what should it be called.

\section{The problem: a belief exposed twice}
\label{sec:problem}

The two exposures have different characters and have been studied by different literatures.

At execution, the question is what a participant reveals in order to trade, and to whom.
This is the territory of market microstructure, where the standard result is that a market
maker able to condition on an informed order will price closer to that order's true value,
so the informed party's surplus falls \cite{glosten1985bidask,kyle1985continuous}. It is
also the territory of maximal extractable value, where an exposed but unexecuted intention
is the entire extraction surface \cite{daian2019flashboys,ethereummev}, and of the order-flow
auctions built to contain it \cite{passeratpalmbach2025dphints}. A request-for-quote
prediction market is structurally an order-flow auction whose order flow is a forecast.

At resolution, the question is who decides what the question meant. When the resolver is a
model, that decision is a computation, and a substantial engineering effort now goes into
proving such computations were performed honestly, through reproducible execution
environments and attested execution inside trusted hardware
\cite{gensyn2026receipts,truebit2026orchestration}. This report argues that those mechanisms
answer a question adjacent to the one that matters, and that the gap between them is where
the interesting residual lives.

I am not aware of work that studies the two exposures together for the same protocol, and
Section~\ref{sec:bg-gap} states that gap on the basis of the literature reviewed there.

\section{A working definition of AI-MEV}
\label{sec:definition}

The definition adopted at the agentic MEV breakout group at MEV-SBC in July 2026 was
deliberately broad: MEV is any information asymmetry leveraged for advantage in a system,
with arbitrage as one instantiation, and the operative question being who obtains systematic
access to capture it \cite{mevsbc2026breakout}. That breadth is useful for framing and
useless for adjudication, because it classifies ordinary market making as MEV. This report
therefore narrows it. The prefix marks the setting rather than adding a criterion: the
traded object is a belief and the resolver is a model, and it is those two facts, not any
clause of the definition, that make the exposures studied here differ from their
counterparts in a human market.

\begin{quote}
\textbf{AI-MEV} is value captured because of \textbf{role-specific protocol access} to
information, timing, execution or resolution power, beyond what is attributable to
forecasting ability or to ordinary competitive compensation for risk borne.
\end{quote}

The qualifier is load-bearing, and it is why Chapter~\ref{ch:evaluation} ends without a
verdict rather than with one. A market maker that prices the information content of an order
it was designed to receive is doing its job. That is adverse selection; it is among the best-established results in market
microstructure \cite{glosten1985bidask}, and calling it extraction would be an error. Any
claim that a mechanism extracts value from informed participants must show that it does
something beyond charging them for being informed.

Four conditions are used in Chapter~\ref{ch:evaluation} to decide whether a measured
transfer qualifies:

\begin{itemize}
  \item \textbf{Asymmetric access.} Some parties see order flow that others do not, or see
    it earlier.
  \item \textbf{Accumulation.} A party observing a large share of flow builds an advantage
    that no single competitive quote compensates.
  \item \textbf{No priced option not to disclose.} The mechanism offers no way to trade
    disclosure off against execution quality.
  \item \textbf{Role-conditional.} The advantage attaches to occupying a protocol position
    rather than to being a better forecaster.
\end{itemize}

\paragraph{How the conditions combine.} A test whose aggregation rule is unstated cannot
adjudicate anything, so the rule is given here. The definition has two limbs: that the
advantage arises from \emph{role-specific protocol access}, and that it exceeds \emph{what
ordinary competitive compensation would provide}. The four conditions are two tests for each
limb. Asymmetric access and role-conditionality bear on the first; accumulation and the
absence of a priced option not to disclose bear on the second. A transfer qualifies as
AI-MEV when at least one condition on each limb holds and the exclusion for risk
compensation does not apply. Neither limb is sufficient alone: an advantage that is
role-specific but fully compensated is a market maker's ordinary return, and an
uncompensated advantage available to anyone is not protocol privilege.

Where no condition holds on either limb, the transfer is ordinary adverse selection and this
report says so. That outcome would be a legitimate result rather than a failure. Stating the
rule also makes it possible for the report to fail to reach a verdict for a locatable
reason, which is what happens in Section~\ref{sec:howmuch}.

One difficulty in the first and fourth conditions should be admitted at the outset, because
Chapter~\ref{ch:evaluation} runs into it and does not resolve it. Both conditions ask
whether access or advantage is protocol-specific. Ethereum's public mempool is open to
anyone who connects, and participation in the MEV that surrounds it is gated by capital and
infrastructure rather than by permission; that is the canonical setting of the literature
this project builds on \cite{daian2019flashboys,ethereummev}. So a reading of
``role-specific protocol access'' strict enough to exclude capital-gated access to a public
stream would exclude most of what MEV means in practice, while a reading permissive enough
to include it makes the conditions weak. This report inherits that tension from the
literature rather than introducing it, applies the conditions as stated, and reports which
verdicts turn on it.

Two terms are used at the point of adjudication and are defined here. \emph{Epistemic reward} is return attributable to
forecasting ability: to knowing something about the world. \emph{Mechanism-mediated
information rent} is return attributable to how the event, the matching rule or the resolver
was specified. In a well-designed market the two coincide. This report is about the
conditions under which they come apart.

\section{Aim and objectives}
\label{sec:aim}

The aim of this project is to establish where strategic informational advantage arises in
AI-native prediction markets, at both the execution and resolution boundaries, to measure
what the execution-side advantage is worth under a controlled counterfactual, and to
determine whether either advantage qualifies as AI-MEV under the definition above rather
than as ordinary adverse selection.

The aim decomposes into five objectives.

\begin{itemize}
  \item \textbf{Establish the mechanisms.} Determine, from primary documentation and
    implementation source rather than from secondary description, what each protocol
    discloses at execution, to whom, and at what moment; and what each delegates at
    resolution, to whom, and under what constraint.
  \item \textbf{Separate the verification properties.} Determine what the verification
    machinery around AI settlement does and does not guarantee, and whether the guarantees
    on offer are substitutes for one another.
  \item \textbf{Measure the execution-side disclosure.} Quantify how much of an informed
    forecaster's expected surplus moves when responders may condition on an observed
    request, against a benchmark identical in every other respect, and establish how that
    quantity varies with the competitive environment and with how much is disclosed.
  \item \textbf{Adjudicate.} Apply the definition and its conditions to the measured
    transfer, and report the verdict, including a verdict of \emph{not established} if that
    is what the evidence supports.
  \item \textbf{Assess mitigations, and bound the claims.} Determine which interventions the
    measurements support, and state throughout what the work demonstrates as against what it
    argues.
\end{itemize}

Objectives one, two, three and five were met. The fourth was not: the adjudication returns
no verdict, for a reason Section~\ref{sec:howmuch} locates precisely. Reporting that as a
failed objective rather than redescribing it as a finding is the honest accounting, and
Chapter~\ref{ch:conclusions} treats it as such.

\section{Contributions}
\label{sec:contributions}

\begin{enumerate}
  \item \textbf{A protocol analysis of two opposite designs}, built from primary
    documentation and implementation source read at a pinned commit, establishing what
    Sapience's relayer broadcasts, to whom, and with what contents, and what Delphi's
    settlement architecture is and has been (Chapter~\ref{ch:protocols}).
  \item \textbf{A separation of three verification properties} that are routinely treated as
    one: semantic validity, reproducibility and execution attestation. The report argues
    that reproducibility and attestation are not substitutes, for a technical reason, and
    that neither protocol's verification machinery guarantees the first. Delphi's own architectural
    history is read as supporting evidence, since a designer with every incentive to close
    the whole problem closed execution discretion and published nothing to indicate that
    specification discretion was closed with it. That is an argument from silence and
    Section~\ref{sec:threelayer-eval} carries the weaker reading of it. This is the report's
    principal contribution and it requires no experiment
    (Chapters~\ref{ch:protocols} and \ref{ch:evaluation}).
  \item \textbf{A measurement of pre-trade disclosure against a matched benchmark.} A
    simulator quantifies how much of an informed forecaster's surplus moves when competing
    responders may condition on the observed request, holding competition, signal quality,
    dependence and risk borne identical. Six parameters are swept
    (Chapters~\ref{ch:method} and \ref{ch:results}).
  \item \textbf{An explicit non-result, reported as one.} Applying the definition above, the
    experiment establishes neither that the measured transfer is AI-MEV nor that it is
    ordinary adverse selection. Under the aggregation rule, the second limb of the definition
    is satisfied and the first is undetermined, because both of the conditions testing it
    turn on the definitional problem set out above; the risk exclusion is set to zero by the
    model rather than tested. What survives without depending on the verdict is a design
    critique: the mechanism offers no priced way to disclose less, at a venue where
    disclosure is complete, address-linked and unpriced, and the entire measured transfer
    accrues to the quoting responders in the model (Chapter~\ref{ch:evaluation}).
  \item \textbf{Two subsidiary measurements with practical consequences.} The share of a
    forecaster's advantage lost to disclosure rises with the forecaster's accuracy, reaching
    about ninety-six per cent at the highest accuracy modelled. And direction alone accounts
    for roughly ninety-five per cent of the transfer, which makes coarsening the disclosed
    size nearly useless as a mitigation and identifies direction as the quantity a mitigation
    would have to protect. Which instrument protects it best is not something this experiment
    can settle (Chapters~\ref{ch:results} and \ref{ch:evaluation}).
\end{enumerate}

\section{Change of direction since the interim report}
\label{sec:intro-change}

The interim report submitted in July 2026 proposed a taxonomy of MEV-like extraction across
decentralised AI networks, supported by two case studies at different layers: a forecasting
subnet and a distributed-training subnet. The final project is narrower. It studies two
prediction-market protocols, one empirical question and one analytical proposition. Three
things drove the change.

The first is that the original breadth would have produced a survey rather than a result. A
taxonomy spanning inference and training, evaluated against no measurement, would have
restated the problem in more categories without settling anything.

The second is that the empirical access the interim assumed proved impractical, which the
interim itself anticipated as an obstacle. Measuring validator behaviour in a live subnet
requires data that is not published. Both of the protocols studied here are documented by
their operators; Delphi's contracts are reported audited by Trail of Bits
\cite{gensyn2026delphi}, and Sapience's implementation is public. Their mechanisms can
therefore be established from primary sources and modelled offline without approaching a
live venue.

The third is the substantive one. The two layers turned out not to be symmetric. In a
distributed-training network the object being evaluated is a model update and the evaluator
is a scoring rule. In an agentic prediction market the object being traded is a belief and
the resolver is itself a model, so the informational asymmetry that MEV describes sits
inside the market's own epistemics rather than beside them. Narrowing to two protocols that
make opposite choices at both boundaries converts a taxonomy into a comparison with a
testable claim.

This is narrowing under evidence, and the material that no longer carries the argument has
been removed rather than retained for its own sake.

\section{Ethical boundary}
\label{sec:ethics-intro}

This project analyses mechanisms in live financial venues, so the boundary it held requires
argument rather than assertion.

No real capital was committed anywhere. No transaction was signed on any network, mainnet or
testnet. No request of any kind was sent to Sapience's relayer, which fans out to real market
makers, so no synthetic order flow entered a real venue. All protocol facts were established
by reading published documentation and open-source implementations; all measurements come
from a simulator.

Chapter~\ref{ch:protocols} nonetheless describes, in detail, an unmitigated information
disclosure in a named protocol that holds real capital. Four things make that defensible.

\begin{itemize}
  \item The broadcast is documented by the protocol's own builder guide, so nothing here is
    disclosed that its operators have not themselves published.
  \item The implementation is public, so the analysis is of published code rather than of
    anything obtained privately or by probing a live system.
  \item No specific market, participant or position is named anywhere in this report.
  \item Mitigations are presented alongside the analysis rather than deferred, in
    Section~\ref{sec:mitigations}.
\end{itemize}

On notification: Sapience was not contacted before this report was written, and the reason
should be stated rather than left to be inferred. Responsible disclosure exists to give an
operator time to remediate something they do not know. The behaviour analysed here is one
the operator documents in its own builder guide and implements in its own open source, so
there was nothing to notify them of that they had not already published. Had the report
found a behaviour undocumented by the protocol, that judgement would have been different. The report contains no usable attack recipe: what it
contributes on this surface is a measurement of what disclosure costs a forecaster and an
argument about how that cost should be classified, not a method for capturing it. On the
classification, Chapter~\ref{ch:evaluation} reaches no verdict, and the ethical position
here does not depend on one. It would be the same if the transfer were shown to be
extraction, because the response this report offers to that possibility is a set of
mitigations rather than a technique.

These commitments are restated in the Declarations.

\section{Report structure}
\label{sec:structure}

Chapter~\ref{ch:background} reviews the five literatures the project draws on and states the
gap between them. Chapter~\ref{ch:protocols} sets Delphi and Sapience side by side, traces
what each exposes at execution and at resolution, and develops the separation between
semantic validity, reproducibility and execution attestation. Chapter~\ref{ch:method}
specifies the disclosure experiment and says before any result what it can and cannot
establish. Chapter~\ref{ch:results} reports the six sweeps. Chapter~\ref{ch:evaluation}
applies the definition above to those results, sets out the limitations at length, and
assesses mitigations. Chapter~\ref{ch:conclusions} concludes.
